\documentclass[12pt,a4paper]{article}
\usepackage[utf8]{inputenc}
\usepackage[T1]{fontenc}
\usepackage[margin=2.5cm]{geometry}
\usepackage{graphicx}
\usepackage{listings}
\usepackage{xcolor}
\usepackage{caption}

\newcommand{\studentgroup}{FAF-242}
\definecolor{termbg}{rgb}{0.96,0.96,0.96}
\definecolor{codegray}{rgb}{0.5,0.5,0.5}

\lstset{
  backgroundcolor=\color{termbg},
  basicstyle=\ttfamily\footnotesize,
  breaklines=true,
  captionpos=b,
  frame=single,
  keepspaces=true,
  numbers=left,
  numbersep=6pt,
  numberstyle=\tiny\color{codegray},
  showstringspaces=false,
  tabsize=2
}

\begin{document}

\begin{titlepage}
    \centering
    \IfFileExists{utm_logo.png}{\includegraphics[width=0.25\textwidth]{utm_logo.png}\\[0.6cm]}{\vspace*{2cm}}
    {\large \textbf{UNIVERSITATEA TEHNICĂ A MOLDOVEI}}\\[0.2cm]
    {\large Facultatea Calculatoare, Informatică și Microelectronică}\\[0.2cm]
    {\large Departamentul Ingineria Software și Automatică}\\[2.5cm]
    {\Large \textbf{LUCRAREA DE LABORATOR NR. 9 \& 10}}\\[0.4cm]
    {\LARGE \textbf{Watching System Calls (Linux) \& Add a System Call (xv6)}}\\[0.4cm]
    {\large Disciplina: Sisteme de Operare}\\[2.5cm]
    \begin{minipage}{0.45\textwidth}\flushleft\textbf{Done by:}\\st. gr. \studentgroup\\Daniel-Ioan \textsc{Strunga}\end{minipage}\hfill
    \begin{minipage}{0.50\textwidth}\flushright\textbf{Checked by:}\\asist. univ. Patricia \textsc{Reițman}\\asist. univ. Marius \textsc{Dumitrașcu}\end{minipage}
    \vfill{\large Chișinău --- 2026}
\end{titlepage}

\newpage
\tableofcontents
\newpage

\section{Scopul lucrării}
Această lucrare combinată are un dublu obiectiv. Pe de o parte (Lab 9 - Track A), se urmărește observarea apelurilor de sistem declanșate de programe reale în Linux folosind utilitarul \texttt{strace} și exersarea apelurilor POSIX fundamentale de I/O (\texttt{open}, \texttt{read}, \texttt{write}, \texttt{close}) omițând bibliotecile de nivel înalt. Pe de altă parte (Lab 10 - Track B), se vizează extinderea nucleului xv6 prin implementarea unui apel de sistem complet nou (\texttt{trace}) capabil să monitorizeze execuția altor syscall-uri pe baza unor măști de biți moștenite între procese.

\section{Laboratorul 9: Watching System Calls (Linux)}
\subsection{Codul sursă POSIX I/O (\texttt{syscalls.c})}
\begin{lstlisting}[language=C]
#include <fcntl.h>
#include <unistd.h>
#include <stdio.h>

int main(void) {
    int fd = open("note.txt", O_WRONLY | O_CREAT | O_TRUNC, 0644);
    if (fd < 0) { perror("open write"); return 1; }
    printf("open returned descriptor %d\n", fd);
    write(fd, "hello, kernel\n", 14);
    printf("write sent 14 bytes\n");
    close(fd);

    fd = open("note.txt", O_RDONLY);
    if (fd < 0) { perror("open read"); return 1; }
    char buf[16];
    read(fd, buf, 14);
    buf[14] = '\0';
    printf("read got 14 bytes: %s", buf);
    close(fd);
    return 0;
}
\end{lstlisting}

\subsection{Ieșirile din terminal (strace)}
\begin{figure}[h!]
\begin{lstlisting}[language=bash]
$ strace -c ls -l >/dev/null
% time     seconds  usecs/call     calls    errors syscall
------ ----------- ----------- --------- --------- ----------------
 17.63    0.000213          12        17           mmap
 13.41    0.000162          10        16           close
  8.36    0.000101           9        11           read
  7.78    0.000094           9        10           openat
  0.83    0.000010          10         1           write
------ ----------- ----------- --------- --------- ----------------
100.00    0.001208                   117        13 total
\end{lstlisting}
\caption{Figure 9.1 -- Statistica syscall-urilor pentru ls -l (117 apeluri totale pe x86\_64)}
\end{figure}

\begin{figure}[h!]
\begin{lstlisting}[language=bash]
$ strace -e trace=openat,write,read,close ./syscalls
openat(AT_FDCWD, "note.txt", O_WRONLY|O_CREAT|O_TRUNC, 0644) = 3
write(1, "open returned descriptor 3\n", 27) = 27
write(3, "hello, kernel\n", 14)         = 14
write(1, "write sent 14 bytes\n", 20)   = 20
close(3)                                = 0
openat(AT_FDCWD, "note.txt", O_RDONLY)  = 3
read(3, "hello, kernel\n", 63)          = 14
write(1, "read got 14 bytes: hello, kernel"..., 33) = 33
close(3)                                = 0
\end{lstlisting}
\caption{Figure 9.2 -- Urmărirea programului C propriu. Se observă cum kernel-ul refolosește descriptorul 3.}
\end{figure}

\subsection{Răspunsuri (Observe Lab 9)}
\begin{enumerate}
    \item \textbf{Cum afișează \texttt{ls} pe ecran?} Comanda \texttt{ls} scrie conținutul pe ecran apelând funcția de sistem \texttt{write}, direcționând datele către descriptorul de fișier \texttt{1} (\texttt{stdout}).
    \item \textbf{De ce fișierul primește descriptorul 3?} La pornirea procesului, kernel-ul asociază implicit 0, 1 și 2. Deoarece regula este alocarea celui mai mic număr liber, primul fișier deschis primește \textbf{3}. După \texttt{close(3)}, numărul redevine disponibil.
    \item \textbf{De ce obținem rezultate diferite fără \texttt{>/dev/null}?} Textul printat de propriul program (ex: \textit{„write sent...”}) va fi interceptat de terminal. Linia conține cuvântul \texttt{write}, așa că regula \texttt{grep} face un match fals, numărând eronat 3 apariții în loc de 2.
\end{enumerate}

\section{Laboratorul 10: Add a System Call: trace (xv6)}
\subsection{Codul sursă user-space (\texttt{user/trace.c})}
\begin{lstlisting}[language=C]
#include "kernel/param.h"
#include "kernel/types.h"
#include "kernel/stat.h"
#include "user/user.h"

int main(int argc, char *argv[]) {
  if (argc < 3) {
    fprintf(2, "Usage: %s mask command\n", argv[0]);
    exit(1);
  }
  if (trace(atoi(argv[1])) < 0) {
    fprintf(2, "%s: trace failed\n", argv[0]);
    exit(1);
  }
  exec(argv[2], &argv[2]);
  exit(0);
}
\end{lstlisting}

\subsection{Ieșirile din terminal (xv6 emulator)}
\begin{figure}[h!]
\begin{lstlisting}[language=bash]
$ trace 32 grep hello README
3: syscall read -> 1023
3: syscall read -> 965
3: syscall read -> 453
3: syscall read -> 0
\end{lstlisting}
\caption{Figure 10.1 -- Izolarea apelurilor de sistem 'read' (masca 32)}
\end{figure}

\begin{figure}[h!]
\begin{lstlisting}[language=bash]
$ trace 2 usertests forkforkfork
usertests starting
test forkforkfork: 5: syscall fork -> 6
6: syscall fork -> 7
7: syscall fork -> 8
...
9: syscall fork -> -1
OK
ALL TESTS PASSED
\end{lstlisting}
\caption{Figure 10.2 -- Verificarea mecanismului de moștenire a măștii în timpul creării de procese multiple}
\end{figure}

\subsection{Răspunsuri (Observe Lab 10)}
\begin{enumerate}
    \item \textbf{De ce 32 izolează \texttt{read}?} Valoarea $32$ este $2^5$. Apelul \texttt{SYS\_read} are ID-ul $5$. Shiftarea la stânga (\texttt{1 << 5}) setează exact bitul corespunzător.
    \item \textbf{Cum se moștenește masca prin \texttt{fork}?} În funcția \texttt{kfork()} am adăugat \texttt{np->trace\_mask = p->trace\_mask;}, procesul copil copiind starea de monitorizare a părintelui.
    \item \textbf{De ce \texttt{echo hi} generează două \texttt{write}?} Programul apelează \texttt{write} o dată pentru șirul \texttt{hi}, și a doua oară pentru terminatorul de rând (\texttt{\textbackslash n}).
    \item \textbf{De ce se loghează \texttt{printk} după apelul sistem?} Pentru a putea prelua și printa valoarea returnată de acel apel, salvată în \texttt{p->trapframe->a0} doar la finalizare.
    \item \textbf{Cum formăm masca pentru \texttt{open} și \texttt{close}?} \texttt{SYS\_open} = 15 ($2^{15} = 32768$), \texttt{SYS\_close} = 21 ($2^{21} = 2097152$). Adunate rezultă \textbf{2129920}.
\end{enumerate}

\section{Concluzii (Comparare Linux - xv6)}
Prin desfășurarea laboratoarelor 9 și 10 am studiat interfața nativă de comunicare dintre aplicațiile din user-space și kernel, implementată printr-un perete de securitate bazat pe syscall-uri. Pe Linux (Lab 9), ne-am bazat pe \texttt{strace} și am observat efortul masiv depus de kernel în fundal și abstractizarea bibliotecii \texttt{libc}. Pe xv6 (Lab 10), am dezvoltat manual un subsistem de audit folosind o mască bitwise performantă ($O(1)$) pentru a izola și urmări apelurile. Ambele medii au confirmat principiul arhitecturii POSIX de bază privind reutilizarea descriptorilor.
\end{document}
